\chapter{Źródła i weryfikacja}
\markboth{Źródła i weryfikacja}{Źródła i weryfikacja}

\section{Co oznacza znaczek \textsuperscript{[?]}}

Znaczek \textsuperscript{\textcolor{accentorange}{[?]}} stawiam przy zdaniu, które zawiera twierdzenie o zdrowiu, nauce, prawie, regulatorze, dacie albo liczbie o świecie, której nie da się wyliczyć w tekście. Znaczy on jedno: autor nie sprawdził tego w źródle pierwotnym w czasie pisania. Nie znaczy, że twierdzenie jest fałszywe, i nie znaczy, że jest prawdziwe. Jest po to, żebyś wiedział, co jest nieprzeczytane.

Przejście znaczka do stanu sprawdzonego wymaga, żeby ktoś otworzył źródło i porównał je ze zdaniem. Samo poczucie, że zdanie brzmi wiarygodnie, nie wystarcza. Liczba oznaczonych miejsc jest w rejestrze stanu w Dodatku Z i jest liczona przez narzędzie, a nie wpisywana.

\section{Jak sprawdzać twierdzenia o śnie}

Nie podaję tu tytułów ani autorów. Autor nie sprawdził ich w tej wersji książki, a wymyślone odnośniki byłyby gorsze niż ich brak. Podaję za to typy źródeł, od najmocniejszych do najsłabszych.

\begin{itemize}
\item \textbf{Wytyczne towarzystw medycyny snu.} Dokumenty, w których zawodowe organizacje podsumowują dowody i formułują zalecenia. Czytaj datę, grupę, dla której napisano zalecenie, i to, czy zalecenie jest mocne, czy warunkowe.
\item \textbf{Przeglądy systematyczne i metaanalizy.} Zbiory badań z opisaną metodą wyszukiwania. Sprawdź, ile badań weszło do przeglądu i jakiej były jakości. Jeden przegląd nie zamyka tematu, a jego wniosek zależy od tego, co do niego trafiło.
\item \textbf{Pojedyncze badania.} Przydatne do formułowania pytania, ale nie do wniosku o sobie. Zwracaj uwagę na wielkość próby, populację badanych i to, czy sen mierzono urządzeniem, czy zgłaszano z pamięci.
\item \textbf{Lekarz, zwłaszcza specjalista medycyny snu.} Jedyne źródło, które zna Twoją historię, wyniki badań i objawy. Żadna książka go nie zastąpi.
\item \textbf{Dane z własnego urządzenia.} Pomocne do obserwacji i porównań w czasie, ale nie są diagnozą. Ich dokładność zależy od urządzenia i od tego, czy producent podał, jak je sprawdzono.
\end{itemize}

Przy każdym twierdzeniu zadaj trzy pytania: kto je podał, na jakich danych i czy ktoś niezależny to powtórzył. Jeśli odpowiedzi brak, traktuj twierdzenie jako hipotezę do sprawdzenia, a nie jako fakt.

\section{Gdy potrzebujesz dwunastu godzin snu stale}

Jeśli Twoja potrzeba snu jest długa od lat, albo jeśli mimo dwunastu godzin jesteś senny, zobacz lekarza, zanim zaplanujesz kolejne tygodnie. Książka nie wie, dlaczego śpisz tyle, i nie powinna zgadywać. Przyczyny długiego snu bywają różne, a część z nich \verify{bywa medyczna i wymaga badania}. Wizyta nie jest porażką planu. Jest jego częścią, którą najłatwiej pominąć.

Przed wizytą zapisz przez kilka tygodni godziny zasypiania i budzenia, to, jak czujesz się w ciągu dnia, czy ktoś zauważył, że chrapiesz, jakie leki i suplementy bierzesz oraz czy sen jest przerywany. Takie notatki są surowcem do rozmowy, nie diagnozą.

\section{Lista kontrolna}

\begin{enumerate}
\item Przy każdym twierdzeniu o zdrowiu znajdź, kto je podał.
\item Sprawdź datę źródła i to, czy nie zastąpiono go nowszym.
\item Odróżnij zalecenie od obserwacji i od wyniku pojedynczego badania.
\item Zauważ, dla jakiej populacji twierdzenie zostało napisane.
\item Nie buduj decyzji na znaczku \textsuperscript{[?]}, dopóki nie otworzysz źródła.
\item Zapisz, co sprawdziłeś i kiedy. Sprawdzenie bez notatki się nie liczy.
\item Jeśli źródło podaje liczbę, policz ją sam z jego danych i porównaj z tym, co zostało podane.
\item Zanim zmienisz sen o więcej niż godzinę, porozmawiaj z lekarzem.
\item Traktuj dane z urządzenia jako tendencję w czasie, a nie jako wynik jednej nocy.
\item Gdy źródła się różnią, zapisz rozbieżność zamiast wybierać wersję, która pasuje do planu.
\end{enumerate}

Ta lista nie gwarantuje, że się nie pomylisz. Zmniejsza natomiast szansę, że pomylisz się bez świadomości.
